\chapter{Pruebas}\label{pruebas}
% tfg-floatbarrier-auto
\makeatletter
\@ifundefined{tfgfloatbarrier}{
  \def\tfgfloatbarrier{1}
  \let\tfgorigfigure\figure
  \let\tfgendorigfigure\endfigure
  \renewenvironment{figure}{\tfgorigfigure}{\tfgendorigfigure\FloatBarrier}
  \let\tfgorigtable\table
  \let\tfgendorigtable\endtable
  \renewenvironment{table}{\tfgorigtable}{\tfgendorigtable\FloatBarrier}
}{}
\makeatother

\FloatBarrier
\section{Introducci\'on}

La fase de pruebas tiene como objetivo validar el sistema desde una doble perspectiva:
\textbf{correctitud funcional} y \textbf{calidad de implementaci\'on}.
En un proyecto con arquitectura desacoplada (frontend y backend),
la validaci\'on debe cubrir tanto el comportamiento interno de los m\'odulos
como los flujos de usuario de extremo a extremo.

Por este motivo se adopta una estrategia de pruebas en capas,
complementada con an\'alisis de mutaci\'on y ejecuci\'on automatizada en CI/CD.

\FloatBarrier
\section{Objetivos de verificaci\'on}

Los objetivos principales de esta fase son verificar que los casos de uso prioritarios funcionan de forma estable, prevenir regresiones al introducir cambios en l\'ogica de negocio, garantizar que las restricciones de seguridad y autorizaci\'on se respetan, medir cobertura y fortaleza de tests de forma objetiva, e integrar la validaci\'on en el ciclo continuo de construcci\'on y despliegue. Esta formulaci\'on conecta calidad funcional y operativa en una misma l\'inea de evidencia.

\FloatBarrier
\section{Estrategia de pruebas en capas}

La estrategia sigue una pir\'amide de testing con cuatro niveles: pruebas unitarias para validar l\'ogica aislada de servicios, utilidades y componentes; pruebas de integraci\'on para validar interacci\'on entre capas y contratos de API; pruebas E2E para validar flujos de usuario en navegador; y pruebas de mutaci\'on para evaluar la capacidad real de los tests de detectar defectos. Esta combinaci\'on equilibra rapidez de feedback y profundidad de verificaci\'on.

La Tabla~\ref{tab:matriz-pruebas} resume las tecnolog\'ias aplicadas.

\begin{table}[!htbp]
\centering
\begin{tabular}{|P{2.9cm}|P{3.6cm}|P{5.4cm}|}
\hline
\textbf{Nivel} & \textbf{Herramientas} & \textbf{Finalidad principal} \\
\hline
Unitarias frontend & Vitest + Testing Library & Validaci\'on de componentes, contextos, servicios y utilidades del cliente. \\
\hline
Unitarias e integraci\'on backend & JUnit + Spring Test + MockMvc & Verificar servicios, controladores, seguridad y persistencia. \\
\hline
E2E frontend & Playwright \cite{playwrightdocs} & Simular escenarios reales de uso en navegador con flujos completos. \\
\hline
Mutaci\'on frontend & Stryker & Medir fortaleza de tests frente a mutantes en servicios y utilidades. \\
\hline
Mutaci\'on backend & PIT \cite{pitest} & Medir capacidad de detecci\'on de fallos en servicios/controladores objetivo. \\
\hline
\end{tabular}
\caption{Matriz de niveles y herramientas de prueba}
\label{tab:matriz-pruebas}
\end{table}

\FloatBarrier
\section{Inventario cuantitativo de pruebas realizadas}

Para documentar de forma expl\'icita el alcance de verificaci\'on,
se ha consolidado el recuento de pruebas por capa a partir de los reportes
generados por Vitest, Playwright y Surefire.

El n\'umero total de pruebas autom\'aticas ejecutadas en la versi\'on actual (MVP final) es:

\[
N_{\text{total}} = N_{\text{frontend unitarias}} + N_{\text{frontend E2E}} + N_{\text{backend}} = 138 + 3 + 850 = 991
\]

El desglose principal es el siguiente: frontend unitarias con 138 casos de prueba en 21 ficheros, frontend E2E con 3 escenarios end-to-end y backend con 850 casos de prueba agregados en 57 suites Surefire. Esta distribuci\'on refleja mayor concentraci\'on en verificaci\'on de servicios.

\FloatBarrier
\section{Pruebas en frontend}

\subsection{Pruebas unitarias y cobertura}

El frontend utiliza Vitest como runner principal,
con entorno \texttt{happy-dom} y configuraci\'on de cobertura por V8.
La suite actual incluye pruebas sobre componentes (navegaci\'on y constructor de estructura ZIP), contextos de autenticaci\'on y tema, p\'aginas clave (login, dashboard, actividades por curso, evaluaciones y perfil), servicios de API (auth, usuarios, cursos, actividades, entregables, entregas, feedback, notificaciones y OneDrive) y utilidades de parsing y validaci\'on de estructura ZIP. Esta cobertura prioriza flujos de uso frecuentes y puntos de integraci\'on con mayor probabilidad de regresi\'on.

Se han definido umbrales m\'inimos de cobertura en la configuraci\'on de Vitest: \texttt{lines} \(\ge 95\%\), \texttt{functions} \(\ge 95\%\), \texttt{statements} \(\ge 95\%\) y \texttt{branches} \(\ge 85\%\). Estos umbrales obligan a mantener un nivel de verificaci\'on consistente.

\subsection{M\'odulos testeados en frontend}

Adem\'as de la cobertura porcentual, se ha realizado un inventario por m\'odulo
de la suite unitaria para dejar evidencia de qu\'e zonas funcionales reciben
mayor esfuerzo de verificaci\'on.

\begin{table}[!htbp]
\centering
\begin{tabular}{|P{5.6cm}|P{2.2cm}|}
\hline
\textbf{M\'odulo frontend} & \textbf{Casos} \\
\hline
Servicios (\texttt{src/services}) & 51 \\
\hline
Componentes (src/components) & 30 \\
\hline
Utilidades (\texttt{src/utils}) & 22 \\
\hline
P\'aginas (\texttt{src/pages}) & 16 \\
\hline
Contextos (\texttt{src/context}) & 11 \\
\hline
Ra\'iz de aplicaci\'on (\texttt{src/}) & 8 \\
\hline
\textbf{Total frontend unitarias} & \textbf{138} \\
\hline
\end{tabular}
\caption{Distribuci\'on de casos unitarios por m\'odulo en frontend}
\label{tab:frontend-modulos-testeados}
\end{table}

\subsection{Pruebas E2E}

Las pruebas end-to-end se ejecutan con Playwright sobre Chromium,
y validan flujos de interacci\'on completos desde la interfaz.

El enfoque de ejecuci\'on adoptado en este proyecto es \textbf{E2E mockeada}:
el navegador se levanta contra el frontend real y se simulan endpoints API
en los escenarios definidos. Esto reduce inestabilidad ambiental y facilita
reproducibilidad en CI.

Se validan, entre otros, los escenarios de autenticaci\'on correcta, autenticaci\'on fallida y flujo de evaluaciones con filtros y descarga mockeada. En la versi\'on actual (MVP final) se ejecutan 3 escenarios y los 3 finalizan en verde; durante la ejecuci\'on pueden aparecer avisos de proxy de Vite para endpoints no mockeados de notificaciones, sin impacto en el resultado de la suite.

\subsection{Mutaci\'on en frontend}

Stryker se aplica sobre c\'odigo de producci\'on de servicios y utilidades,
excluyendo archivos de test.
El objetivo es verificar que los tests no solo ejecutan l\'ineas,
sino que detectan alteraciones sem\'anticas en el comportamiento.

La cobertura mediante mutaci\'on en frontend se ha realizado con un procedimiento que selecciona el alcance en \path{src/services/**/*.ts} y \path{src/utils/**/*.ts}, excluye expl\'icitamente tests y archivos de infraestructura (\path{*.test.*}, \path{*.spec.*}, \path{api.ts}, \path{index.ts}), ejecuta \texttt{npm run test:mutation} generando mutantes y relanzando Vitest por lotes, y analiza mutantes \textit{survived}/\textit{timeout} para reforzar casos de prueba d\'ebiles. Este enfoque prioriza secciones con mayor impacto funcional.

En la versi\'on actual (MVP final), \texttt{npm run test:mutation} instrumenta 10 ficheros de servicios y utilidades con 435 mutantes. El resultado global es un \textit{mutation score} del 92.10\%, con 267 mutantes eliminados, 23 supervivientes y 1 \textit{timeout}. Como criterio de aceptaci\'on de esta fase, Stryker aplica umbrales \texttt{high=75}, \texttt{low=60} y \texttt{break=50}, siendo este \'ultimo el que determina fallo de pipeline.

\FloatBarrier
\section{Pruebas en backend}

\subsection{Pruebas unitarias e integraci\'on}

En backend se emplea JUnit junto con Spring Test y MockMvc.
La validaci\'on cubre l\'ogica de servicios de dominio, validaciones y reglas de negocio, contratos REST de controladores, seguridad y autorizaci\'on por contexto, repositorios sobre H2, entidades y conversiones DTO. Este alcance garantiza consistencia entre reglas de negocio, persistencia y API expuesta.

Entre los tests de referencia de mayor carga funcional destaca
\texttt{EntregaServiceTest}, que valida escenarios de entrega,
restricciones y evoluci\'on de estado.

\subsection{M\'odulos testeados en backend}

El backend se ha cuantificado a partir de nodos \texttt{testcase} en reportes generados por Maven Surefire (el plugin oficial de construcci\'on y ejecuci\'on de pruebas de Java que orquesta la fase de test en el ciclo de vida del proyecto). Esto permite reflejar correctamente pruebas parametrizadas y escenarios anidados.

\begin{table}[!htbp]
\centering
\begin{tabular}{|P{5.6cm}|P{2.2cm}|}
\hline
\textbf{M\'odulo backend} & \textbf{Casos} \\
\hline
Servicios (\texttt{service}) & 542 \\
\hline
Controladores (\texttt{controller}) & 149 \\
\hline
Repositorios (\texttt{repository}) & 47 \\
\hline
DTOs y validaciones (\texttt{dto}) & 40 \\
\hline
Seguridad (\texttt{security}) & 31 \\
\hline
Entidades (\texttt{entity}) & 22 \\
\hline
Manejo de excepciones (exception) & 10 \\
\hline
Configuraci\'on (\texttt{config}) & 8 \\
\hline
Aplicaci\'on base (\texttt{root}) & 1 \\
\hline
\textbf{Total backend} & \textbf{850} \\
\hline
\end{tabular}
\caption{Distribuci\'on de casos de prueba por m\'odulo en backend}
\label{tab:backend-modulos-testeados}
\end{table}

\subsection{Cobertura backend}

La cobertura se genera con JaCoCo \cite{jacoco} (Java Code Coverage Library, un marco de trabajo de an\'alisis de c\'odigo abierto para medir c\'omo el c\'odigo fuente es ejercitado por los tests autom\'aticos) durante la fase de test, produciendo informe HTML en
\texttt{backend/target/site/jacoco/index.html}. En la versi\'on actual (MVP final), el reporte agregado indica 89.55\% de instrucciones, 70.31\% de ramas, 90.75\% de l\'ineas y 91.60\% de m\'etodos cubiertos. Este reporte se utiliza como soporte de revisi\'on t\'ecnica para detectar m\'odulos subprobados y orientar ampliaciones de suite.

\subsection{Mutaci\'on en backend}

El perfil Maven \texttt{mutation} ejecuta PIT (PIT Mutation Testing, un sistema de pruebas de mutaci\'on para la m\'aquina virtual de Java que altera artificialmente el bytecode de los programas para medir qu\'e porcentaje de esos defectos es capaz de detectar la suite de pruebas) sobre clases objetivo de servicios y controladores seleccionados.
La configuraci\'on define umbrales de mutaci\'on y cobertura,
y genera reportes estructurados para su an\'alisis.

En esta implementaci\'on se mutan de forma controlada 10 clases objetivo:
7 de servicio y 3 de controlador. La selecci\'on cubre integraci\'on con
OneDrive, validaci\'on ZIP, feedback, mapeo de entidades, contexto de
seguridad, recuperaci\'on de contrase\~na, Cloudinary y los controladores
asociados a OneDrive, OAuth de Microsoft y feedback. Esta selecci\'on
concentra el esfuerzo de mutaci\'on en servicios de integraci\'on y flujo
cr\'itico.

El flujo seguido para cobertura de mutaci\'on en backend consiste en ejecutar el perfil con \texttt{mvn -Pmutation ... mutationCoverage}, generar mutantes con operadores de retorno, aritm\'etica e incrementos, relacionar mutante y test para identificar qu\'e prueba mata cada defecto sint\'etico, y revisar mutantes supervivientes para crear pruebas adicionales o ajustar aserciones. Este proceso mejora la eficacia real de la suite.

En la versi\'on actual (MVP final), PIT genera 144 mutantes, de los que 134 son eliminados por la suite. El resultado es un \textit{mutation score} del 93\%, una cobertura de l\'inea del 95\% sobre clases mutadas y una fuerza de test del 96\%. El perfil exige adem\'as un m\'inimo de mutaci\'on del 70\% y de cobertura del 75\%, reforzando que la cobertura no sea solo de l\'ineas, sino de detecci\'on efectiva de fallos.

\FloatBarrier
\section{Automatizaci\'on en CI/CD}

La validaci\'on se integra en pipelines de GitHub Actions,
lo que garantiza ejecuci\'on repetible por cada cambio relevante.

Workflows principales: \texttt{frontend-ci.yml} para lint, unitarias con cobertura, E2E y build de frontend; \texttt{full-stack-ci.yml} para verificaci\'on backend (\texttt{mvn verify}) y frontend completo; y \texttt{mutation-ci.yml} para mutaci\'on backend (PIT) y frontend (Stryker) con subida de artefactos. Esta organizaci\'on separa capas de validaci\'on sin perder coherencia global.

Esta automatizaci\'on reduce el riesgo de integraci\'on tard\'ia, y permite detectar regresiones antes de integrar c\'odigo en la rama principal y desplegarlo en producci\'on.

\FloatBarrier
\section{Resultados verificados}

De acuerdo con el informe de estrategia de testing del proyecto
(corte de verificaci\'on final), se obtienen los siguientes resultados representativos:

\begin{table}[!htbp]
\centering
\begin{tabular}{|P{4.4cm}|P{7.4cm}|}
\hline
\textbf{M\'etrica} & \textbf{Resultado} \\
\hline
E2E frontend (Playwright) & 3 escenarios ejecutados, 3 en verde. \\
\hline
Cobertura frontend & Statements 97.78\%, Branches 90.68\%, Functions 98.20\%, Lines 98.25\%. \\
\hline
Suite frontend & 138 tests aprobados en 21 archivos durante la ejecuci\'on del comando test:coverage. \\
\hline
Suite backend completa & 850 casos en 57 suites Surefire, 0 fallos y 0 errores. \\
\hline
Total pruebas autom\'aticas & 991 ejecuciones (138 unitarias frontend + 3 E2E + 850 backend). \\
\hline
Mutaci\'on frontend (Stryker) & Score global 92.10\% sobre 435 mutantes (por encima del umbral \texttt{break}=50). \\
\hline
Mutaci\'on backend (PIT) & 144 mutantes generados, mutation score 93\%, line coverage 95\%, test strength 96\%. \\
\hline
Test objetivo backend & \texttt{EntregaServiceTest}: 89 tests, 0 fallos, 0 errores. \\
\hline
\end{tabular}
\caption{Resumen de resultados de verificaci\'on}
\label{tab:resultados-pruebas}
\end{table}

Estos valores muestran una madurez de testing consistente con el nivel de exigencia
de un TFG de orientaci\'on profesional.

\FloatBarrier
\section{Criterios de aceptaci\'on}

Como criterio operativo de calidad se establecen m\'inimos de \texttt{lint} sin errores, suites unitarias, integraci\'on y E2E en verde, coberturas por encima de umbrales configurados, m\'etricas de mutaci\'on por encima de umbral de ruptura y ausencia de regresiones funcionales en flujos cr\'iticos. Estos criterios permiten cierre objetivo de iteraciones.

\FloatBarrier
\section{Limitaciones y riesgos}

Aunque la estrategia es robusta, se identifican l\'ineas de mejora como ampliar E2E con backend real en entorno de staging para ciertos casos, extender mutaci\'on a mayor superficie de clases backend y reforzar pruebas de carga y comportamiento en condiciones extremas. Estas mejoras consolidar\'ian la robustez en escenarios de producci\'on.

No obstante, las limitaciones actuales no comprometen la validez de la verificaci\'on
para el alcance definido del proyecto.

\FloatBarrier
\section{Conclusiones}

La estrategia de pruebas aplicada combina cobertura, profundidad y automatizaci\'on,
aportando evidencia objetiva de calidad funcional y t\'ecnica.

La incorporaci\'on de mutaci\'on como criterio adicional permite evaluar
la eficacia real de la suite, superando una visi\'on meramente cuantitativa
de cobertura.

En conjunto, los resultados respaldan que el sistema dispone de una base
de validaci\'on s\'olida para mantenimiento, evoluci\'on y despliegue continuo.






